\section{Reading a majority by differentiation}\label{sec:vote}

Suppose each active site in a group \(C\) holds a symbol and a strict
majority hold the same symbol \(y\). We want to send \(\pi(y)\) to a
destination \(U\), for a projection \(\pi\) specified after the
resources for the gate have been chosen. The cost per output bit will
be the largest price of a wire from one voter to \(U\). There is also
a fixed allowance for all the possible symbols. On unanimity the
gate returns almost this entire allowance in a single fixed tensor.

The construction keeps one slot for each possible symbol. Each slot
carries a tensor supported at a fixed total degree. A voter subtracts
a fixed amount of degree in the slot for its own symbol. We choose
that amount so a strict majority exhausts an ordinary slot. A
derivative doubles the available degree in exactly one slot, and only
the term in which that slot belongs to the winner survives.

\subsection{Degree tensors and the vote}

For a nonempty set \(C\) of sites and an integer \(m\ge1\), put
\[
 W_m(C)=
 \sum_{\substack{e_i\ge0\ (i\in C)\\ \sum_{i\in C}e_i=m-1}}
         \bigotimes_{i\in C}v_{e_i},
\]
where \(v_e\) is the basis vector indexed by degree \(e\).
Equivalently, \(W_m(C)\) takes the coefficient of \(t^{m-1}\) in a
product in \(\mathbb F[t]/(t^m)\). If \(m\) is a power of the
characteristic, changing the basis at every site to
\((1+t)^x\), \(0\le x<m\), gives the constraint
\(\sum_i x_i\equiv m-1\pmod m\). Indeed, \((1+t)^m=1\) in this
algebra, and among the basis polynomials only \((1+t)^{m-1}\) has a
term of degree \(m-1\).

We use the spectral point from the compatible Euclidean prices of
Theorem~\ref{thm:flat}. For these prices \(W_m(C)\) has the same value
as a common label of size \(m\) when \(m\) is a power of the
characteristic. The estimate at other sizes will measure what the gate
returns. Section~\ref{sec:degree} proves both statements.

\begin{lemma}[Price of a fixed total degree]\label{lem:degree-width}
Over an algebraically closed field of characteristic \(p>0\), for the
compatible Euclidean prices and every nonempty \(C\),
\[
              h(W_m(C))=w(C)\log m
              \qquad(m\text{ a power of }p).
\]
For any integer \(s\ge1\) and \(|C|\ge2\), one has
\[
              h(W_s(C))\ge
              w(C)\log\left\lceil\frac{s}{|C|-1}\right\rceil .
\]
In particular \(h(W_s(C))\ge w(C)\log s-O(\log(2|C|))\) for all
nonempty \(C\), with an absolute implied constant.
\end{lemma}

\begin{lemma}[A vote with one wire]\label{lem:vote}
Let \(C\) be a nonempty set of sites and \(U\) a site outside \(C\).
Each voter holds a \(b\)-bit symbol and an activity flag. The domain
fixes the number \(n\ge1\) of active voters. There is an allowance
\(j\ge0\) such that for each \(0\le l\le b\) a gate taking a public
projection \(\pi\colon\{0,1\}^b\to\{0,1\}^l\) as a control has cost at most
\[
               j+l\max_{i\in C}w(\{i,U\})+\log |C|.
\]
When a strict majority of the active voters hold \(y\), the gate
appends \(e_{\pi(y)}\) at \(U\) and a nonzero tensor \(J\) on \(C\).
For each such input the same \(J\) works for every choice of \(l,\pi\).
When all active voters hold \(y\), a permutation of symbol slots,
determined by \(y\) and applied at every site of \(C\), takes \(J\) to
a fixed tensor of price at least \(j-O(\log(2|C|))\).
\end{lemma}

\begin{proof}
Choose a power \(m\) of \(p\) with \(m\ge100(n+1)^2\). Put
\(M=2^b\) and \(j=Mw(C)\log m\). Use slots with degree basis
\(v_0,\ldots,v_{2m}\), padding smaller tensors with zeros.

\emph{Doubling the degree.}
For \(z\in\mathbb F\), define the entry of \(V(z)\) at a degree tuple
\((e_i)\) to be the coefficient of \(t^{m-1}\) in the remainder of
\(t^{\sum_i e_i}\) modulo \(t^m-zt\). Reduction of a monomial of
degree at least \(m\) lowers its degree by \(m-1\) and multiplies it
by \(z\). At total degree \(a>0\) divisible by \(m-1\), the entry is
\(z^{a/(m-1)-1}\). All other entries vanish. Thus \(V(0)=W_m(C)\),
and \(V'(0)=W_{2m-1}(C)\) has total degree \(2(m-1)\).

At \(z\ne0\) the polynomial \(t^m-zt\) has \(m\) distinct roots,
since its derivative is \(-z\). Interpolating at those roots expresses
the coefficient of \(t^{m-1}\) as a weighted sum of \(m\) evaluations.
Each evaluation of a product factors across the sites. Consequently
\(V(z)\) is a restriction of \(G_C(m)\), of price at most \(w(C)\log m\).

\emph{Sending the projection.}
Index the \(M\) slots by symbols. Define the polynomial tensor
\begin{equation}\label{eq:vote-derivative}
 T_\pi(z)=\sum_{x\in\{0,1\}^b} e_{\pi(x)}^{(U)}
           \otimes\bigotimes_{u\in\{0,1\}^b}
             \begin{cases} V'(z),&u=x,\\ V(z),&u\ne x. \end{cases}
\end{equation}
At \(z\ne0\), entrywise differentiation multiplies the entry of \(V(z)\)
by \((\sum_i e_i-(m-1))/((m-1)z)\). The denominator is invertible
because \(p\) divides \(m\). This multiplier is a sum of local
diagonal operators. At voter \(i\) use \(e_i/((m-1)z)\), subtracting
\(1/z\) at one designated voter.

For the summand corresponding to voter \(i\), supply \(V(z)^{\otimes M}\)
at cost at most \(j\) and share an \(l\)-bit label \(v\) between \(i\)
and \(U\). At \(i\), apply the sum of its local operators on the slots
\(x\) with \(\pi(x)=v\), and erase \(v\). Site \(U\) keeps \(v\).
Summing over the voters appends \(T_\pi(z)\) for every choice of the
public control \(\pi\). A summand pays for one wire, and summing
\(|C|\) gates adds at most \(\log|C|\) to the largest cost by
\eqref{eq:sum-prices}. For each input tensor, including its public
controls and any auxiliary factors, this output is polynomial in \(z\).
Specialization at zero therefore satisfies the same cost bound.

\emph{Selecting the winner.}
Write \(a_0=\lfloor n/2\rfloor+1\) and
\(d=\lfloor(m-1)/a_0\rfloor+1\).
Each active voter locally replaces \(v_e\) by \(v_{e-d}\) in the slot
for its symbol, sending \(v_e\) to zero when \(e<d\). Let \(t_x\)
count the votes for \(x\). If \(t_y\ge a_0\), then \(t_yd>m-1\).
Every term in \eqref{eq:vote-derivative} at zero vanishes except the
one with the derivative in slot \(y\). For \(x\ne y\),
\(t_xd\le(a_0-1)d\le m-1\), since \(m-1\ge a_0(a_0-1)\).
For the winning slot, put \(s=2(m-1)-nd+1\). From \(2a_0-n\ge1\)
and \(d\le(m-1)/a_0+1\) we get
\[
             s\ge(m-1)/a_0-n+1\ge m/(2(n+1)).
\]
There is thus enough degree in the winning slot even if all \(n\)
voters choose it.

In the surviving term, each tuple of remaining degrees lifts uniquely
by adding \(d\) at voters for the corresponding symbol. The lifted
degrees sum to the prescribed degree and are at most \(2(m-1)\).
It follows that \(J\) is a product of fixed-degree tensors, with
\(W_{\,2(m-1)-t_y d+1}(C)\) in slot \(y\) and \(W_{\,m-t_x d}(C)\) in
each other slot \(x\). All factors are nonzero. This description also
shows that \(J\) is independent of \(l,\pi\).

On unanimity, moving slot \(y\) to the front gives the tensor
\(W_s(C)\otimes W_m(C)^{\otimes(M-1)}\), which depends only on the fixed
parameters. By Lemma~\ref{lem:degree-width} its price is at least
\[
     j-w(C)\log(m/s)-O(\log(2|C|))
                         \ \ge\ j-O(\log(2|C|)),
\]
since \(n\le |C|\) and \(w(C)\le1\).
\end{proof}

\subsection{Why the degree tensor has the price of a shared label}\label{sec:degree}

Lemma~\ref{lem:vote} needs the price of \(W_m(C)\) for the resource it
supplies and the price of \(W_s(C)\) for the resource it returns. For
general \(s\), a direct degeneration extracts a smaller common label.
To obtain the exact price at powers of \(p\), we let the sites vary in
Euclidean space. The partial diagonals and the degree tensors have
the same price on a segment. A degeneration on intersecting sets of
sites will force their prices to agree on every convex hull.

For these comparisons use the joint Euclidean assignment of
\cite[Theorem~\ext{G}{thm:main} and Proposition~\ext{G}{prop:joint-convexity}]{oaG}.
Its unnormalized price for a nonzero tensor is jointly convex in
the placements of its factors, invariant under Euclidean isometries
and homogeneous under simultaneous dilation. At coincident sites the
factors form a single local space. Adding a scalar factor does not
change the price, and placing all factors together gives price zero.
On a fixed finite set \(S\) of positive width the unnormalized price
is \(D h\), where
\begin{equation}\label{eq:euclidean-price}
 D=\delta(S),\qquad
 w(P)=\frac{\delta(P)}{D},\qquad
 \delta(P)=\sqrt{2\pi}\,\mathbb E\max_{v\in P}\langle g,v\rangle ,
 \qquad g\sim N(0,I).
\end{equation}
The quantity \(\delta\) depends only on the convex hull, adds under
Minkowski sums and equals length on segments.

\begin{proof}[Proof of Lemma~\ref{lem:degree-width}]
Fix a power \(m>1\) of \(p\). For a nonempty finite set \(A\) in the
Euclidean flat, let \(f(A)\) be the unnormalized price of \(W_m(A)\)
divided by \(\log m\). The remainder construction of Section~\ref{sec:vote} gives a restriction
of \(G_A(m)\) at each \(z\ne0\) and gives \(W_m(A)\) at zero. Polynomial
specialization therefore gives \(0\le f(A)\le\delta(A)\). To prove equality we adapt the mean-width
argument for diagonal tensors in \cite[proof of Theorem~\ext{G}{thm:flatlaw}]{oaG}.

\emph{Intersecting sets of sites.}
For finite sets \(A,B\) with nonempty intersection \(I=A\cap B\), we
claim that
\begin{equation}\label{eq:degree-submod}
                 f(A\cup B)+f(I)\le f(A)+f(B).
\end{equation}
Represent \(W_m(A)\otimes W_m(B)\) by coefficient extraction at
\(t^{m-1}u^{m-1}\) in \(\mathbb F[t,u]/(t^m,u^m)\). The local
inputs are polynomials in \(t\) on \(A\setminus I\), in \(u\) on
\(B\setminus I\), and in both variables on \(I\). In the inputs from
\(B\setminus I\), substitute \(t+zu\) for \(u\).

At \(z\ne0\) this changes the form only by invertible local maps and a
nonzero scalar. Indeed, \(t\mapsto t,\ u\mapsto t+zu\) is a graded
automorphism of the quotient algebra because
\((t+zu)^m=t^m+z^m u^m\). Apply it to the full product after applying
its inverse to the inputs at \(I\). The automorphism scales the
one-dimensional top-degree component by \(z^{m-1}\).
At \(z=0\) every site contributes to the \(t\)-degree, and precisely
the sites in \(I\) contribute to the \(u\)-degree. The specialized
form is \(W_m(A\cup B)\otimes W_m(I)\). Taking prices proves
\eqref{eq:degree-submod}.

\emph{From segments to polytopes.}
Restricting the local degree at a new site to zero recovers the old
tensor, so \(f\) is monotone under inclusion. At a repeated site the
map \(v_e\mapsto\sum_{a+b=e} v_a\otimes v_b\) inserts an extra slot,
and selecting degree zero in that slot reverses the map. If a new
site lies in \(\operatorname{conv}(A)\), convexity bounds its price
by the price of adjoining it at one of the old sites. Monotonicity
gives the reverse bound. Hence \(f(A)\) depends only on
\(K=\operatorname{conv}(A)\).

The resulting function on polytopes satisfies
\(f(K+L)\le f(K)+f(L)\). Take vertex lists \((x_i)\) of \(K\) and
\((y_j)\) of \(L\), and give \(W_m\) one slot for each pair \((i,j)\).
For the placements \(x_i\) and \(y_j\) the values are \(f(K)\) and
\(f(L)\), since repeated sites have no effect. At the midpoint
placement the convex hull is \((K+L)/2\), so joint convexity and
homogeneity give the inequality.

For a vector \(v\) and an integer \(N\ge1\), form the chain of
translates \(A,A+v,\ldots,A+Nv\). When adjoining the adjacent pair
\((A+hv)\cup(A+(h+1)v)\) to the preceding union, the intersection
already contains \(A+hv\). By \eqref{eq:degree-submod} and translation
invariance, each increment in price is at most \(f(K+[0,v])-f(K)\).
The whole union contains \(x,x+Nv\) for any \(x\in A\). At two sites
\(W_m\) is a diagonal after reversal of one degree basis. Since its
unnormalized price divided by \(\log m\) is the distance, we obtain
\[
 N\|v\|\le f(K+[0,Nv])
              \le f(K)+N\bigl(f(K+[0,v])-f(K)\bigr).
\]
Dividing by \(N\) and taking \(N\to\infty\) gives
\(f(K+[0,v])\ge f(K)+\|v\|\). Subadditivity gives the other inequality.
Iterating from a point shows that \(f=\delta\) on every finite sum
of segments.

\emph{Averaging rotations.}
If two polytopes \(K,L\) have Hausdorff distance at most \(t\), each
lies in the sum of the other with \(tQ\), where \(Q\) is a fixed cube
containing the unit ball. Monotonicity and subadditivity give
\(\lvert f(K)-f(L)\rvert\le t\delta(Q)\). Thus \(f\) extends
continuously to compact convex sets. The Minkowski average of all
orthogonal images of \(K\) is a ball \(B\) centered at the origin,
since its support function is constant on the unit sphere.
Approximating the rotation average by finite averages gives
\(f(B)\le f(K)\) by subadditivity and invariance. The same averaging
applied to a segment expresses any origin-centered ball as a limit of
finite Minkowski sums of segments. Hence \(f(B)=\delta(B)\).
Invariance and Minkowski additivity of \(\delta\) give
\(\delta(B)=\delta(K)\), so
\[
                   \delta(K)=f(B)\le f(K)\le\delta(K).
\]
Normalizing as in \eqref{eq:euclidean-price} proves the power-of-\(p\)
formula. At \(m=1\) the tensor is the scalar unit.

\emph{Arbitrary remaining degree.}
Let \(|C|=a+1\ge2\). At \(a\) sites of \(W_s(C)\) retain degrees
\(x_1,\ldots,x_a\) in \(\{0,\ldots,\lfloor(s-1)/a\rfloor\}\).
At the final site retain degree \(s-1-a x_0\) with \(x_0\) in the
same interval. The remaining support satisfies \(\sum_{i=1}^a x_i=a x_0\).
Scale each term by \(z\) to the power
\[
      \sum_{i=1}^a x_i^2-a x_0^2=\sum_{i=1}^a(x_i-x_0)^2 .
\]
On the left the exponent is a sum of local integer weights, so the
scaling is local for \(z\ne0\). On the right it is nonnegative on
the support and vanishes precisely at equality. Specializing at zero
therefore gives a diagonal of size \(1+\lfloor(s-1)/a\rfloor
=\lceil s/a\rceil\), whose price is the claimed lower bound. For a
single site both the price and \(w(C)\) are zero.
\end{proof}

